% ---
% id: foundations/req:tiers
% kind: req
% title: "Reporting of the tiers"
% summary: >-
%   Every run that reports a key rate reports it at each tier it reaches and
%   names the strongest, whose number is the headline and gives the key length.
%   No module labels a number with a tier its checks do not support, and a
%   request for an unknown tier is refused.
% prominence: standard
% status: accepted
% uses: [def:tiers]
% mentions: [meat-core/frank-wolfe/req:certified-evaluation]
% assumes: []
% refs: []
% provenance: newly-derived
% lean:
% ---

A run reaches the \tref{certified} tier only where its certified evaluation
succeeds, which is the business of the solver that forms the number
(\xref{meat-core/frank-wolfe/req:certified-evaluation}); the rule below only
says how the tiers a run reaches are reported.

\begin{requirement}\label{req:tiers}
Every run that reports a key rate reports it at every tier of
Definition~\ref{def:tiers} that it reaches: the estimate always, the
first-order number wherever the hypotheses of the claim that establishes it
held, and the certified number wherever the requirement of the solver that
forms it reports it at that tier, which may withhold the tier from a number
that meets Definition~\ref{def:tiers}\,(iii) but rests on an upstream
premise; together with the field
\iface{reported_tier}, which names the strongest tier reached, the tier that
carries the run's claim. Where the first-order number is reported, the number
at that strongest tier is the headline key rate, and the key length the run
reports is formed from it; a number at a weaker tier is recorded beside it.
No module reports a number under a tier stronger than the checks that number
has passed support, and a request naming a tier that
Definition~\ref{def:tiers} does not list is refused as an unknown tier.
\end{requirement}
